\documentclass[border=5mm]{standalone}
% ==============================================================================
% SETUP.TEX - CONFIGURACIÓN GLOBAL DEL PROYECTO
% ==============================================================================
% Este archivo centraliza todos los paquetes y estilos.
% Debe incluirse en el preámbulo del libro principal y de los archivos standalone.
% \PassOptionsToPackage{gray}{xcolor}
% --- 1. CODIFICACIÓN E IDIOMA ---
% fix-cm habilita las fuentes Computer Modern en tamaños arbitrarios (por debajo
% de 5 pt, entre otros). Debe cargarse antes que fontenc.
\usepackage{fix-cm}
\usepackage[utf8]{inputenc}
\usepackage[T1]{fontenc}
\usepackage[spanish, es-tabla]{babel}  
\usepackage{csquotes} % Recomendado para babel y citas

\usepackage{import}
% mode=image: Usa el PDF pre-compilado, nunca intenta recompilar.
% Debes compilar cada figura manualmente antes de compilar el libro.
\usepackage[mode=image, subpreambles=false]{standalone}

% --- 2. MATEMÁTICAS Y CIENCIAS ---
\usepackage{amsmath, amssymb, amsfonts, mathtools}
\usepackage{enumitem}

% LaTeX trae \DeclareMathSizes{5}{5}{5}{5}, así que a 5 pt (\tiny) los subíndices
% salen del mismo tamaño que la base: en las etiquetas de figuras el M de
% $\omega_M$ queda desproporcionado. Se restablece la proporción habitual.
\DeclareMathSizes{5}{5}{3.7}{3}

% --- 3. DISEÑO Y GEOMETRÍA --- 
% Unificamos los márgenes para todo el documento
\usepackage[left=2.5cm,right=2.5cm,top=3cm,bottom=3cm]{geometry}
\makeatletter
\ifstandalone % Evita que geometry dañe el recorte de las figuras bajándolas 3cm
  \@ifclasswith{standalone}{crop=false}{
    % Es un capítulo/sección: Dejamos que geometry actúe.
  }{
    % Es una figura: Neutralizamos geometry para que no las recorte mal.
    \geometry{pass}
  }
\fi
\makeatother
\usepackage{parskip} % Espaciado entre párrafos sin sangría
\usepackage{float}   % Para usar [H] en figura

% Ajustes de flotantes para evitar espacios verticales exagerados
\renewcommand{\topfraction}{0.95}      % Permite que las figuras ocupen hasta el 95% superior
\renewcommand{\bottomfraction}{0.95}   % Permite que las figuras ocupen hasta el 95% inferior
\renewcommand{\textfraction}{0.05}     % Permite hojas con tan solo un 5% de texto
\renewcommand{\floatpagefraction}{0.8} % Requiere que una página de solo figuras esté 80% llena
\setcounter{topnumber}{4}
\setcounter{bottomnumber}{4}
\setcounter{totalnumber}{10}

% --- 4. GRÁFICOS Y TABLAS ---
\usepackage{graphicx}
\usepackage{booktabs} % Tablas profesionales
\usepackage{caption}  % Control de captions y \captionof
\usepackage{tikz}
\usetikzlibrary{arrows.meta, calc, angles, quotes, babel, backgrounds}
\usepackage{pgfplots}
\usepgfplotslibrary{groupplots}
\usepgfplotslibrary{external}
\usepgfplotslibrary{patchplots}
\pgfplotsset{compat=1.18}
\usepackage[american, straightvoltages]{circuitikz}

% --- 5. COLORES Y CAJAS ---

\usepackage[table]{xcolor}
\usepackage{tcolorbox}

% Definición de colores corporativos/estilo
\definecolor{utnblue}{RGB}{0, 51, 102}
\definecolor{codegreen}{rgb}{0,0.6,0}
\definecolor{codegray}{rgb}{0.5,0.5,0.5}
\definecolor{codepurple}{rgb}{0.58,0,0.82}
\definecolor{backcolour}{rgb}{0.96,0.96,0.96}

% --- 6. ENTORNOS PERSONALIZADOS (CAJAS) ---

% Caja "Resultado" (Azul Corporativo) — Definiciones, fórmulas clave, resultados importantes.
% Uso: \begin{resultado}[Fórmula de Euler] ... \end{resultado}
\newtcolorbox{resultado}[1][]{
    colback=utnblue!5!white,
    colframe=utnblue,
    title=\textbf{#1},
    fonttitle=\bfseries,
    sharp corners,
    boxrule=0.5mm
}

% Caja "Nota" (Gris) — Advertencias, aplicaciones, contexto secundario.
% Uso: \begin{nota}[Cuidado con la calculadora] ... \end{nota}
% Uso: \begin{nota} ... \end{nota}  → título por defecto "Nota"
\newtcolorbox{nota}[1][Nota]{
    colback=codegray!5!white,
    colframe=codegray,
    title=\textbf{#1},
    fonttitle=\bfseries,
    sharp corners,
    boxrule=0.5mm
}

% Caja inline para resaltar un valor y enlazarlo con su otra aparición en una tabla
% o en una cuenta: mismo color, mismo valor.
% Uso: \marcavalor{$4$}  o  \marcavalor[codegreen]{$4$}
\newtcbox{\marcavalor}[1][utnblue]{
    on line,
    boxsep=1pt, left=2pt, right=2pt, top=1pt, bottom=1pt,
    colback=#1!10!white,
    colframe=#1,
    boxrule=0.4pt,
    arc=2pt
}

% --- 7. CÓDIGO FUENTE (LISTINGS) ---
\usepackage{listings}

\lstdefinestyle{miestilo}{
    backgroundcolor=\color{backcolour},   
    commentstyle=\color{codegreen},
    keywordstyle=\color{magenta},
    numberstyle=\tiny\color{codegray},
    stringstyle=\color{codepurple},
    basicstyle=\ttfamily\footnotesize,
    breakatwhitespace=false,         
    breaklines=true,                 
    captionpos=b,                    
    keepspaces=true,                 
    numbers=left,                    
    numbersep=5pt,                  
    showspaces=false,                
    showstringspaces=false,
    showtabs=false,                  
    tabsize=2,
    frame=single,                    
    rulecolor=\color{codegray}
}

\lstset{style=miestilo}
% Soporte para tildes y ñ en bloques de código
\lstset{literate={ñ}{{\~n}}1 {á}{{\'a}}1 {é}{{\'e}}1 {í}{{\'i}}1 {ó}{{\'o}}1 {ú}{{\'u}}1}

% Operadores matemáticos personalizados

% Vector en sentido discreto (lista finita de coordenadas): negrita + flecha.
% Uso: \vect{v}, \vect{e}_k, \vect{0}.
\newcommand{\vect}[1]{\vec{\mathbf{#1}}}

% Fecha de versión y comando para capítulos standalone
\providecommand{\verdate}{\the\year.\ifnum\month<10 0\fi\the\month.\ifnum\day<10 0\fi\the\day}
\newcommand{\chapterversion}{%
    \vspace{-1.5cm}\begin{flushright}\small\textit{\color{codegray}Versión~\verdate}\end{flushright}\vspace{0.5cm}%
}

% --- 8. HIPERVÍNCULOS (DEBE SER EL ÚLTIMO PAQUETE) ---
\usepackage[colorlinks=true, linkcolor=utnblue, urlcolor=utnblue, citecolor=utnblue]{hyperref}

% --- 9. REFERENCIAS CRUZADAS ENTRE CAPÍTULOS STANDALONE ---
% Cuando se compila un capítulo standalone, xr-hyper lee los .aux de los
% otros capítulos (si existen) para resolver \ref{} a labels externos.
% En la compilación del libro completo este bloque no se activa.
\ifstandalone
  \usepackage{xr-hyper}
  \makeatletter
  \newcommand{\xrchap}[1]{%
    \edef\xrc@self{\detokenize{Capitulo#1}}%
    \edef\xrc@job{\jobname}%
    \ifx\xrc@self\xrc@job\else
      \IfFileExists{../#1capitulo/Capitulo#1.aux}%
        {\externaldocument{../#1capitulo/Capitulo#1}}{}%
    \fi
  }
  \makeatother
  \xrchap{01}\xrchap{02}\xrchap{03}\xrchap{04}
  \xrchap{05}\xrchap{06}\xrchap{07}\xrchap{08}
  \xrchap{09}\xrchap{10}\xrchap{11}\xrchap{12}
  \xrchap{13}
\fi
\usetikzlibrary{shadows.blur}

\begin{document}
    \begin{tikzpicture}[>=Latex, font=\small,
        block/.style={draw=utnblue!60, line width=0.5pt, fill=utnblue!8,
                      minimum width=2.6cm, minimum height=1.1cm, rounded corners=2pt,
                      blur shadow={shadow blur steps=4, shadow blur extra rounding=0pt,
                                   shadow xshift=0.4pt, shadow yshift=-0.6pt, shadow opacity=25}},
        signal/.style={->, thick, utnblue},
        digisignal/.style={->, thick, red!70!black},
        label_style/.style={font=\footnotesize}
    ]

        % --- Señal analógica de entrada ---
        \coordinate (in) at (0,0);

        % Bloque ADC
        \node[block, fill=orange!12, minimum width=1.6cm] (adc) at (2.4,0) {\textbf{ADC}};

        % Bloque Procesamiento
        \node[block, fill=red!8] (proc) at (7,0) {\parbox{2.2cm}{\centering\textbf{Procesamiento}\\\textbf{Digital}}};

        % Bloque DAC
        \node[block, fill=orange!12, minimum width=1.6cm] (dac) at (11.6,0) {\textbf{DAC}};

        % Salida
        \coordinate (out) at (14,0);

        % --- Conexiones ---
        \draw[signal] (in) -- (adc.west) node[label_style, above, midway] {$x(t)$};
        \draw[digisignal] (adc.east) -- (proc.west) node[label_style, above, midway] {$x[n]$};
        \draw[digisignal] (proc.east) -- (dac.west) node[label_style, above, midway] {$y[n]$};
        \draw[signal] (dac.east) -- (out) node[label_style, above, midway] {$y(t)$};

        % Mini ondas decorativas
        % Onda analógica de entrada con ruido Fuerte
        \begin{scope}[shift={(0.8,0.95)}]
            \draw[utnblue, thick] plot[domain=-1:1, samples=200] (\x, {-0.25*sin(deg(\x*pi)) + 0.06*rand});
            \node[label_style, utnblue] at (0, 0.65) {Analógica};
        \end{scope}

        % Stems digitales de x[n] (Muestreo de onda ruidosa pero un poco desplazada arriba)
        \begin{scope}[shift={(4.5,0.95)}]
            \foreach \x/\h in {-1.0/0.05, -0.8/0.11, -0.6/0.30, -0.4/0.18, -0.2/0.22, 0.0/-0.08, 0.2/-0.10, 0.4/-0.28, 0.6/-0.18, 0.8/-0.22, 1.0/0.05} {
                \draw[red!70!black, thick] (\x, 0) -- (\x, \h);
                \fill[red!70!black] (\x, \h) circle (1pt);
            }
            \node[label_style, red!70!black] at (0, 0.65) {Digital ($x[n]$)};
            \draw[thin, black!40] (-1.1,0) -- (1.1,0);
        \end{scope}

        % Stems digitales de y[n] (Senoidal limpia, ya procesada, sin ruido)
        \begin{scope}[shift={(9.5,0.95)}]
            \foreach \x/\h in {-1.0/0.0, -0.8/0.15, -0.6/0.24, -0.4/0.24, -0.2/0.15, 0.0/0.0, 0.2/-0.15, 0.4/-0.24, 0.6/-0.24, 0.8/-0.15, 1.0/0.0} {
                \draw[red!70!black, thick] (\x, 0) -- (\x, \h);
                \fill[red!70!black] (\x, \h) circle (1pt);
            }
            \node[label_style, red!70!black] at (0, 0.65) {Digital ($y[n]$)};
            \draw[thin, black!40] (-1.1,0) -- (1.1,0);
        \end{scope}

        % Onda continua de salida (Sin ruido, senoidal pura reconvertida)
        \begin{scope}[shift={(13.2,0.95)}]
            \draw[utnblue, thick] (-1.0, 0) sin (-0.5, 0.25) cos (0, 0) sin (0.5, -0.25) cos (1.0, 0);
            \node[label_style, utnblue] at (0, 0.65) {Analógica};
        \end{scope}

    \end{tikzpicture}
\end{document}
